%% Chiral Chevalley-Eilenberg factorisation algebra: generators and
%% relations, complementary to celestial_holography.tex
%% subsec:costello-e3-chiral-ce-gen-rel. Fully explicit Feynman-traced
%% generators/relations for the holomorphic factorisation algebra
%% on chiral CE (co)chains, with Costello-Francis-Gwilliam 2026
%% first-principles rigor. Vol II, HT-landscape / 3d-gravity bridge.

%% Local macro safety (providecommand only)
\providecommand{\SCchtop}{\mathsf{SC}^{\mathrm{ch},\mathrm{top}}}
\providecommand{\hCS}{\mathrm{hCS}}
\providecommand{\Obs}{\mathsf{Obs}}
\providecommand{\cA}{\mathcal{A}}
\providecommand{\cB}{\mathcal{B}}
\providecommand{\cE}{\mathcal{E}}
\providecommand{\cF}{\mathcal{F}}
\providecommand{\cL}{\mathcal{L}}
\providecommand{\cM}{\mathcal{M}}
\providecommand{\cO}{\mathcal{O}}
\providecommand{\cV}{\mathcal{V}}
\providecommand{\cW}{\mathcal{W}}
\providecommand{\cG}{\mathcal{G}}
\providecommand{\cH}{\mathcal{H}}
\providecommand{\cP}{\mathcal{P}}
\providecommand{\cR}{\mathcal{R}}
\providecommand{\cU}{\mathcal{U}}
\providecommand{\cD}{\mathcal{D}}
\providecommand{\fg}{\mathfrak{g}}
\providecommand{\fh}{\mathfrak{h}}
\providecommand{\fsl}{\mathfrak{sl}}
\providecommand{\fgl}{\mathfrak{gl}}
\providecommand{\dbar}{\bar\partial}
\providecommand{\CE}{\mathrm{CE}}
\providecommand{\MC}{\mathrm{MC}}
\providecommand{\Sym}{\mathrm{Sym}}
\providecommand{\id}{\mathrm{id}}
\providecommand{\Ran}{\mathrm{Ran}}
\providecommand{\Tr}{\mathrm{Tr}}
\providecommand{\Fact}{\mathrm{Fact}}
\providecommand{\End}{\mathrm{End}}
\providecommand{\Conf}{\mathrm{Conf}}
\providecommand{\Disk}{\mathrm{Disk}}
\providecommand{\Ainf}{A_\infty}
\providecommand{\Linf}{L_\infty}
\providecommand{\Fobs}{\mathbf{Obs}}
\providecommand{\Omegabar}{\Omega^{0,\bullet}}

\chapter{Chiral Chevalley--Eilenberg Factorisation Algebra:
Generators, Relations, and Feynman Traces}
\label{ch:chiral-ce-factalg-gen-rel}

%%%%%%%%%%%%%%%%%%%%%%%%%%%%%%%%%%%%%%%%%%%%%%%%%%%%%%%%%%%%%%%%%%%%%%%%%%%%%
\section*{Prologue}
%%%%%%%%%%%%%%%%%%%%%%%%%%%%%%%%%%%%%%%%%%%%%%%%%%%%%%%%%%%%%%%%%%%%%%%%%%%%%

The chiral Chevalley--Eilenberg cochain algebra
$C^\bullet_{\mathrm{ch}}(\fg,M)$ on a Riemann surface $C$ is the
specific factorisation algebra through which the Costello--Gwilliam BV
quantisation of holomorphic-topological gauge theory lands in the
$\infty$-category of $E_1$-chiral algebras. On the Koszul locus, its
factorisation homology over a transverse disc recovers the universal
enveloping of the iterated loop algebra $\fg\otimes\cO_C$; Dolbeault
Fubini degenerates the $E_3$-structure on $Y=\R^3\times C\times\C^2$
along $C$, and the conformal vector of the non-critical Sugawara
construction converts the chiral $E_1$ direction into the third
$E_1$-factor of the $E_3$-algebra on the transverse $3$-disc; the
chain-level topologisation witness is
Theorem~\ref{thm:ch-ce-topologization-E3-top} below, the
$(\infty,1)$-categorical promotion is
Theorem~\ref{thm:chiral-higher-deligne}\,(2) of
Chapter~\ref{ch:chiral-higher-deligne}.

The subsection
(\S\ref{subsec:costello-e3-chiral-ce-gen-rel}) of
\texttt{celestial\_holography.tex} establishes the skeleton: the
graded vector space, the $\Disk_3$-action, and the leading
Feynman diagrams. The chapter that follows inscribes the full
generators-and-relations presentation: explicit tridegree accounting,
named propagators (Bergman kernel, Green function, prime form),
Feynman-traced $n$-loop OPE coefficients with motivic-MZV provenance,
Wilson-line factorisation modules, and the Ayala--Francis
factorisation-homology Fubini descent from $\Disk_3$ to
$\Disk_1$. Every coefficient in every OPE is tracked to its
$\zeta(k)/(2\pi i)^k$ motivic representative; every structural map is
identified with a named Costello--Gwilliam renormalisation-scale
propagator.

Seven propositions and six theorems crystallise the chapter:
the chain-level $C^\bullet_{\mathrm{ch}}(\fg,M)$ as a
$\Disk_3$-algebra
(Theorem~\ref{thm:ch-ce-disk3-algebra});
the factorisation-homology Fubini from $\Disk_3$ to $\Disk_1$ via
Ayala--Francis--Tanaka
(Theorem~\ref{thm:ch-ce-fubini-disk3-to-disk1});
the explicit Weiss-cover descent on a pair of discs
(Theorem~\ref{thm:ch-ce-weiss-descent});
the one-loop OPE with motivic-$\zeta(2)$ coefficient
(Theorem~\ref{thm:ch-ce-oneloop-motivic});
the $\hbar^2 = -1/8$ two-loop sunrise for $\fsl_2$ via a named
Feynman integral
(Theorem~\ref{thm:ch-ce-sunrise-sl2-minus-eighth});
the Wilson-line factorisation module as path-ordered exponential
(Theorem~\ref{thm:ch-ce-wilson-line-module});
and the 6d hCS identification on $Y=\R^3\times C\times\C^2$ specialising
to the K3 avatar on $C=E$ with 24 punctures
(Theorem~\ref{thm:ch-ce-k3-avatar-24-punctures}).

%%%%%%%%%%%%%%%%%%%%%%%%%%%%%%%%%%%%%%%%%%%%%%%%%%%%%%%%%%%%%%%%%%%%%%%%%%%%%
\section{The graded vector space and the chiral CE differential}
\label{sec:ch-ce-graded-vs-differential}
%%%%%%%%%%%%%%%%%%%%%%%%%%%%%%%%%%%%%%%%%%%%%%%%%%%%%%%%%%%%%%%%%%%%%%%%%%%%%

\begin{definition}[Chiral Chevalley--Eilenberg cochains on a Riemann
surface]
\label{def:ch-ce-cochains}
\index{chiral Chevalley--Eilenberg cochains}%
Fix a smooth complex curve $C$, a finite-dimensional Lie algebra
$\fg$ with invariant pairing $\langle-,-\rangle\colon\fg\otimes\fg\to\C$,
and a $\fg$-module $M$. The \emph{chiral Chevalley--Eilenberg cochain
complex on $C$ with coefficients in $M$} is the tridegreed complex
\begin{equation}
\label{eq:ch-ce-cochains-def}
C^\bullet_{\mathrm{ch}}(\fg, M)(C)
\;:=\;
\Omega^{0,\bullet}(C)
\;\otimes\;
\Sym^\bullet(\fg^\vee[-1])
\;\otimes\;
M,
\end{equation}
with tridegree $(p,q,r)$ where $p\in\{0,1\}$ is the Dolbeault degree,
$q\in\Z_{\ge 0}$ is the CE polynomial degree in $\fg^\vee[-1]$, and
$r$ is the internal $M$-degree. Total cohomological degree is
$p+q+r$. The differential is
\begin{equation}
\label{eq:d-ch}
d_{\mathrm{ch}}
\;=\;
\dbar_C \;+\; d_{\CE},
\end{equation}
with $\dbar_C\colon\Omega^{0,p}\to\Omega^{0,p+1}$ the Dolbeault
operator on $C$ and $d_{\CE}$ the Chevalley--Eilenberg differential
on $\Sym^\bullet(\fg^\vee[-1])\otimes M$, which on a generating
cochain $\xi\in\fg^\vee[-1]$ is
\begin{equation}
\label{eq:d-CE-generator}
d_{\CE}\xi
\;=\;
\tfrac12\xi([-,-])
\;+\;
\rho_M(\xi)\cdot(-),
\end{equation}
with $\rho_M\colon\fg\to\End(M)$ the module action.
\end{definition}

\begin{proposition}[Tensor-orthogonal nilpotency]
\label{prop:ch-ce-nilpotency}
\index{chiral CE differential!nilpotency}%
\ClaimStatusProvedHere
The differential $d_{\mathrm{ch}}$ of \eqref{eq:d-ch} satisfies
$d_{\mathrm{ch}}^2 = 0$.
\end{proposition}

\begin{proof}
Expand
\begin{equation}
d_{\mathrm{ch}}^2
\;=\;
\dbar_C^2 + \{\dbar_C, d_{\CE}\} + d_{\CE}^2.
\end{equation}
Each summand vanishes separately. First, $\dbar_C^2 = 0$ is the Dolbeault
integrability on a smooth curve, i.e.\ $\dbar\circ\dbar = 0$ on
$\Omega^{0,\bullet}(C)$. Second, the anticommutator
$\{\dbar_C, d_{\CE}\}$ vanishes on the nose: $\dbar_C$ acts on the
$\Omega^{0,\bullet}(C)$ tensor factor, $d_{\CE}$ acts on the
$\Sym^\bullet(\fg^\vee[-1])\otimes M$ tensor factor, and the two factor
algebras commute as graded operators on the tensor product
(tensor-orthogonality). Third, $d_{\CE}^2 = 0$ is the classical
Chevalley--Eilenberg identity, equivalent to the Jacobi identity on
$\fg$ together with the module axiom on $M$. Each identity is
independently verified: $\dbar^2=0$ by direct computation in
local coordinates ($\dbar^2 f = \partial_{\bar z_1}\partial_{\bar z_2}f
- \partial_{\bar z_2}\partial_{\bar z_1}f = 0$); $d_{\CE}^2 = 0$ by the
Jacobi-Koszul sign rule verified on generating $\xi_a,\xi_b,\xi_c
\in\fg^\vee[-1]$. Tensor-orthogonality follows from the functorial
definition of the tensor product of chain complexes.
\end{proof}

\begin{remark}[Primary provenance]
\label{rem:ch-ce-primary}
The complex $C^\bullet_{\mathrm{ch}}(\fg,M)$ is the chain-level
realisation of Francis--Gaitsgory's chiral-Koszul dual of
$U(\fg_C)$~\cite[Thm.\,6.4.2]{FG12}: on a formal disc,
$C^\bullet_{\mathrm{ch}}(\fg,\C)$ computes $\mathrm{R}\mathrm{Hom}_{U(\fg\otimes\cO)}(\C,\C)$,
the derived endomorphism of the trivial module over the loop algebra.
\end{remark}

\begin{definition}[Dolbeault tridegree conventions]
\label{def:tridegree-conventions}
\index{tridegree conventions}%
On the formal disc $D_z\subset C$ centred at $z$ with local coordinate
$w$, the Dolbeault degree $p$ counts the number of $d\bar w$-factors,
the CE degree $q$ counts the polynomial degree in the linear ghost
coordinates $\xi^a\in\fg^\vee[-1]$, and the $M$-degree $r$ is internal.
Total ghost number $\mathrm{gh}(x) = p+q+r$ is additive under tensor
product. The $z$-weight (holomorphic conformal weight) $\Delta(x)$
counts the number of $\partial_w$-derivatives, and the $\bar z$-weight
$\bar\Delta(x)$ counts the number of $\partial_{\bar w}$-derivatives
applied to the component functions. All four gradings are preserved
by $\dbar_C$ (up to $+1$ in $p$) and by $d_{\CE}$ (up to $+1$ in $q$).
\end{definition}

%%%%%%%%%%%%%%%%%%%%%%%%%%%%%%%%%%%%%%%%%%%%%%%%%%%%%%%%%%%%%%%%%%%%%%%%%%%%%
\section{Explicit generators on a formal disc}
\label{sec:ch-ce-generators}
%%%%%%%%%%%%%%%%%%%%%%%%%%%%%%%%%%%%%%%%%%%%%%%%%%%%%%%%%%%%%%%%%%%%%%%%%%%%%

\begin{definition}[Elementary generators of the chiral CE algebra]
\label{def:ch-ce-generators}
\index{chiral CE generators}%
Fix a basis $\{T_a\}_{a=1}^{\dim\fg}$ of $\fg$ orthonormal under
$\langle-,-\rangle$, with dual basis $\{T_a^\vee\}$ of $\fg^\vee$ and
dual generating cochains $\{\xi^a\}\subset\fg^\vee[-1]$. Over a formal
disc $D_z\subset C$ with coordinate $w$, the generators of
$C^\bullet_{\mathrm{ch}}(\fg,M)(D_z)$ as a commutative dg
$\Omegabar(D_z)$-algebra are:
\begin{enumerate}[label=\textup{(G\arabic*)},ref=G\arabic*]
\item\label{gen:ch-A}
$\cA^a(w) \;=\; A^a_{\bar w}(w)\,d\bar w \;\in\;
\Omega^{0,1}(D_z)\otimes\fg$, tridegree $(1,0,0)$, $z$-weight $0$,
$\bar z$-weight $0$;
\item\label{gen:ch-B}
$\cB^a(w) \;=\; B^a(w)\,\xi^a \;\in\;
\Omega^{0,0}(D_z)\otimes\fg^\vee[-1]$, tridegree $(0,1,0)$,
$z$-weight $0$;
\item\label{gen:ch-dA}
$\partial^k\cA^a(w) = \partial_w^k A^a_{\bar w}(w)\,d\bar w$
for $k\ge 1$, tridegree $(1,0,0)$, $z$-weight $k$;
\item\label{gen:ch-dB}
$\partial^k\cB^a(w) = \partial_w^k B^a(w)\,\xi^a$ for $k\ge 1$,
tridegree $(0,1,0)$, $z$-weight $k$;
\item\label{gen:ch-M}
module generators $m_\alpha(w)\in\Omega^{0,0}(D_z)\otimes M$
(with $\alpha$ indexing a basis of $M$), tridegree $(0,0,r_\alpha)$
where $r_\alpha$ is the internal $M$-degree of the basis element.
\end{enumerate}
Normal-ordered products are the Wick-ordered bilinears
\begin{equation}
\label{eq:normal-ordered-products}
\normord{\cA^a\cA^b}(z)
\;:=\;
\lim_{w\to z}
\Bigl[\cA^a(w)\cA^b(z) - \frac{\langle T^a,T^b\rangle}{(w-z)^2}
\Bigr],
\end{equation}
with the subtraction removing the double pole. Analogous definitions
apply to $\normord{\cA^a\cB^b}$ and $\normord{\cB^a\cB^b}$.
\end{definition}

\begin{proposition}[Strict $\Omega^{0,\bullet}$-linearity and Leibniz]
\label{prop:ch-ce-omega-linearity}
\ClaimStatusProvedHere
The generators \eqref{gen:ch-A}--\eqref{gen:ch-M} of
Definition~\ref{def:ch-ce-generators} satisfy:
\begin{enumerate}[label=\textup{(\roman*)}]
\item $\dbar_C(f\cdot\cA^a) = (\dbar f)\cdot\cA^a$ for every
$f\in\Omega^{0,0}(D_z)$;
\item $d_{\CE}(\cA^a\cB^b) = (d_{\CE}\cA^a)\cB^b + (-)^{|\cA^a|}
\cA^a(d_{\CE}\cB^b)$;
\item $d_{\CE}\cA^a = 0$ (the gauge field is CE-closed), and
$d_{\CE}\cB^a = \tfrac12 f^a_{bc}\cB^b\cB^c$
(the antighost carries the dual of the Lie bracket).
\end{enumerate}
\end{proposition}

\begin{proof}
Clauses~(i) and~(ii) are restatements of the graded Leibniz rule in
each tensor factor; $\dbar_C$ is an $\Omega^{0,0}$-derivation on
$\Omega^{0,1}$ because $\dbar(f\cdot g\,d\bar w) = (\dbar f)g\,d\bar w
\pm f(\dbar g)\,d\bar w$, and the $\dbar g$ term is a $d\bar w$-form
wedged with $d\bar w$ which vanishes on a one-dimensional (complex)
curve. Clause~(iii): $\cA^a$ lives in tensor factor
$\Omega^{0,1}\otimes\fg$ with no $\fg^\vee$-content, so
$d_{\CE}\cA^a = 0$. The antighost $\cB^a = B^a(w)\xi^a$ has
$d_{\CE}\xi^a = \tfrac12\xi^a([-,-]) = \tfrac12 f^a_{bc}\xi^b\xi^c$
(using $\xi^a([T_b,T_c]) = \xi^a(f^d_{bc}T_d) = f^a_{bc}$ by
orthonormality), whence
$d_{\CE}\cB^a = B^a(w)\cdot\tfrac12 f^a_{bc}\xi^b\xi^c = \tfrac12
f^a_{bc}B^a(w)\xi^b\xi^c$; renaming summation indices gives the stated
form $\tfrac12 f^a_{bc}\cB^b\cB^c$ when $B^a$ is absorbed into
$\cB^a$.
\end{proof}

\begin{remark}[Matching BV-BRST labelling of 6d hCS]
\label{rem:ch-ce-match-6d-bv}
\index{BV-BRST!correspondence to chiral CE}%
The generator \eqref{gen:ch-A} corresponds to the 6d hCS field
$A\in\Omega^{0,1}(X,\mathrm{ad})$ of~\cite[Def.~5.1]{CFG2026} restricted
along the $C$-direction; \eqref{gen:ch-B} corresponds to the antighost
$c^+\in\Omega^{0,3}$ in the hCS BV-BRST complex
(Definition~\ref{def:hcs-bv-brst} of
Chapter~\ref{ch:six-d-hcs-e3-chiral-avatar-platonic}) after Serre-dual
pairing with the K3 direction. The derivative generators
\eqref{gen:ch-dA}--\eqref{gen:ch-dB} correspond to the Taylor modes of
$A$ and $c^+$ around a point of $C$; they span the dual of the
jet-bundle $J^\infty(\fg\otimes\cO_C)$ of the trivial bundle.
\end{remark}

%%%%%%%%%%%%%%%%%%%%%%%%%%%%%%%%%%%%%%%%%%%%%%%%%%%%%%%%%%%%%%%%%%%%%%%%%%%%%
\section{Factorisation axioms and Weiss-cover descent}
\label{sec:ch-ce-factorisation-axioms}
%%%%%%%%%%%%%%%%%%%%%%%%%%%%%%%%%%%%%%%%%%%%%%%%%%%%%%%%%%%%%%%%%%%%%%%%%%%%%

\begin{definition}[Factorisation algebra on $C$]
\label{def:ch-ce-factalg}
\index{factorisation algebra!chiral CE}%
The assignment
\begin{equation}
\label{eq:factalg-defn}
U \;\mapsto\; \cF_{\fg,M}(U)
\;:=\;
\mathrm{R}\Gamma\!\bigl(U,\;C^\bullet_{\mathrm{ch}}(\fg,M)|_U\bigr)
\;=\;
\bigl(\Omega^{0,\bullet}(U)\otimes\Sym^\bullet(\fg^\vee[-1])\otimes M,\;
d_{\mathrm{ch}}\bigr)
\end{equation}
from open subsets $U\subseteq C$ to cochain complexes, with structure
maps described below, is the \emph{chiral CE factorisation algebra
on $C$}.
\end{definition}

\begin{theorem}[Tensor-product axiom on disjoint unions]
\label{thm:ch-ce-tensor-axiom}
\index{factorisation algebra!tensor-product axiom}%
\ClaimStatusProvedHere
For disjoint opens $U_1,\ldots,U_n\subseteq C$, the canonical map
\begin{equation}
\label{eq:tensor-factor-iso}
\cF_{\fg,M}(U_1)\otimes\cdots\otimes\cF_{\fg,M}(U_n)
\;\longrightarrow\;
\cF_{\fg,M}(U_1\sqcup\cdots\sqcup U_n)
\end{equation}
is a cochain quasi-isomorphism. In particular, for two disjoint
discs $D_1\sqcup D_2$,
\begin{equation}
\cF_{\fg,M}(D_1\sqcup D_2)
\;\simeq\;
\cF_{\fg,M}(D_1)\otimes\cF_{\fg,M}(D_2).
\end{equation}
\end{theorem}

\begin{proof}
The cochain complex factors through two tensor sides. The Dolbeault
factor $\Omega^{0,\bullet}$ satisfies the tensor axiom by direct
inspection: $\Omega^{0,\bullet}(U_1\sqcup U_2) =
\Omega^{0,\bullet}(U_1)\oplus\Omega^{0,\bullet}(U_2)$ as graded vector
spaces (disjoint support), and $\dbar$ acts block-diagonally. The
Sym-factor $\Sym^\bullet(\fg^\vee[-1])$ is a finite-dimensional
polynomial algebra per point and is constant on $U$ (pulled back along
$U\to\mathrm{pt}$), so
$\Sym^\bullet(\fg^\vee[-1])\otimes M$ contributes a constant tensor
factor that is manifestly tensor-multiplicative on disjoint unions. The
$d_{\mathrm{ch}}$ differential is the sum $\dbar_C + d_{\CE}$ with the
two pieces acting on the two separated factors, so the
quasi-isomorphism is strict.
\end{proof}

\begin{theorem}[Structure map for nested opens]
\label{thm:ch-ce-weiss-descent}
\index{factorisation algebra!Weiss-cover descent}%
\ClaimStatusProvedHere
For $U_1\subset V$, $U_2\subset V$ mutually disjoint and both
contained in a larger open $V\subseteq C$, the factorisation structure
map
\begin{equation}
\label{eq:structure-map-nested}
m_{U_1,U_2;V}\colon
\cF_{\fg,M}(U_1)\otimes\cF_{\fg,M}(U_2)
\;\longrightarrow\;
\cF_{\fg,M}(V)
\end{equation}
is given by (a) extending each $\Omega^{0,\bullet}(U_i)$-section by zero
to $V\setminus U_i$ and summing as sections over $V$, (b) taking the
commutative-algebra product on the $\Sym^\bullet(\fg^\vee[-1])\otimes
M$ tensor factor, and (c) composing with the $d_{\mathrm{ch}}$-acyclic
descent map from
$\Omega^{0,\bullet}(U_1)\oplus\Omega^{0,\bullet}(U_2)$ to
$\Omega^{0,\bullet}(V)$.

The structure map is strictly associative and symmetric, and
Weiss-cover descent holds: for a Weiss cover $\{U_i\}_{i\in I}$ of $V$
(a cover by disjoint unions of discs with each finite subset of $V$
contained in some $U_i$),
\begin{equation}
\label{eq:weiss-descent}
\cF_{\fg,M}(V)
\;=\;
\mathrm{hocolim}_{i\in I}\cF_{\fg,M}(U_i),
\end{equation}
the homotopy colimit in the $\infty$-category of cochain complexes.
\end{theorem}

\begin{proof}
Associativity of $m_{U_1,U_2;V}$ is a direct consequence of the
commutativity of $\Omega^{0,\bullet}$-multiplication and of the
$\Sym^\bullet$-multiplication. Symmetry: swapping $U_1\leftrightarrow
U_2$ produces a sign $(-)^{|x_1||x_2|}$ from Koszul rule on the
$\Omega^{0,\bullet}$ factor and another $(-)^{|y_1||y_2|}$ from the
$\Sym^\bullet$ factor; the product of the two signs is
$(-)^{|x_1||x_2|+|y_1||y_2|}$, which matches the graded-commutativity
of the tensor product of two commutative dg algebras.

Weiss-cover descent follows from the Dolbeault Mayer--Vietoris sequence
on $V$: covering by discs $\{D_i\}$, the $\mathrm{\check{C}ech}$ complex
$\check C^\bullet(\{D_i\},\Omega^{0,\bullet}|_V)$ computes
$\Omega^{0,\bullet}(V)$ up to quasi-isomorphism
(Costello--Gwilliam~\cite[Prop.\,4.3.4]{CG17}). The
$\Sym^\bullet(\fg^\vee[-1])\otimes M$ factor is a constant sheaf and
contributes the identity under Weiss descent. The homotopy colimit
\eqref{eq:weiss-descent} is therefore the standard Dolbeault
sheaf-cohomology of $\Omega^{0,\bullet}(V)$ tensored with the constant
factor.
\end{proof}

\begin{proposition}[Explicit nilpotency on two-disc overlap]
\label{prop:ch-ce-two-disc-nilpotency}
\index{factorisation algebra!structure-map nilpotency}%
\ClaimStatusProvedHere
For two disjoint discs $D_1 = \{|w-z_1|<r_1\}$ and $D_2 = \{|w-z_2|<r_2\}$
with $|z_1-z_2|>r_1+r_2$, contained in $V = \{|w|<R\}$, the
structure-map composite
\begin{equation}
\label{eq:two-disc-composite}
d_{\mathrm{ch}}^V \circ m_{D_1,D_2;V}
\;=\;
m_{D_1,D_2;V}\circ\bigl(d_{\mathrm{ch}}^{D_1}\otimes\id + \id\otimes
d_{\mathrm{ch}}^{D_2}\bigr)
\end{equation}
is a nilpotent chain map of total degree $+1$. In particular,
$d_{\mathrm{ch}}^V$ applied to the structure-map image
$m_{D_1,D_2;V}(x_1\otimes x_2)$ satisfies
\begin{equation}
d_{\mathrm{ch}}^V m_{D_1,D_2;V}(x_1\otimes x_2)
\;=\;
m_{D_1,D_2;V}(d_{\mathrm{ch}}^{D_1}x_1\otimes x_2 + (-)^{|x_1|}
x_1\otimes d_{\mathrm{ch}}^{D_2}x_2).
\end{equation}
\end{proposition}

\begin{proof}
The structure map $m_{D_1,D_2;V}$ is an $\Omega^{0,\bullet}$-linear map
over the common disc of definition (after extension by zero on
$V\setminus(D_1\sqcup D_2)$). The Dolbeault $\dbar_V$ and
$d_{\CE}^V$ act on disjoint tensor factors, and extension-by-zero
commutes with $\dbar_V$ because the Dolbeault operator is local: on an
open neighbourhood of $\partial D_1\cap V$, the section
$\iota^V_{D_1}(x_1)$ vanishes identically, so $\dbar_V\iota^V_{D_1}(x_1)$
also vanishes outside $D_1$, and on $D_1$ the operators agree by local
inspection.
\end{proof}

\begin{remark}[Primary-literature provenance]
\label{rem:ch-ce-factalg-primary}
Theorem~\ref{thm:ch-ce-tensor-axiom} matches
\cite[Def.\,3.1.4]{CostelloGwilliam2017} applied to the free commutative
dg algebra on a two-term linear factorisation algebra (Dolbeault +
ghost); Theorem~\ref{thm:ch-ce-weiss-descent} matches
\cite[Prop.\,4.3.4]{CostelloGwilliam2017} with the Weiss-cover
descent conditions specialised to a Riemann surface.
Proposition~\ref{prop:ch-ce-two-disc-nilpotency} is a chain-level
witness of the prefactorisation axiom at $n=2$.
\end{remark}

%%%%%%%%%%%%%%%%%%%%%%%%%%%%%%%%%%%%%%%%%%%%%%%%%%%%%%%%%%%%%%%%%%%%%%%%%%%%%
\section{The Disk-3 algebra action and factorisation-homology Fubini}
\label{sec:ch-ce-disk3-fubini}
%%%%%%%%%%%%%%%%%%%%%%%%%%%%%%%%%%%%%%%%%%%%%%%%%%%%%%%%%%%%%%%%%%%%%%%%%%%%%

\begin{theorem}[$\Disk_3$-algebra on $Y = \R^3\times C\times\C^2$]
\label{thm:ch-ce-disk3-algebra}
\index{$\Disk_3$-algebra!chiral CE}%
\ClaimStatusProvedElsewhere\,\cite[Thm.\,4.4.1]{CG17}
Let $Y = \R^3\times C\times\C^2$ with the product holomorphic
volume form
$\Omega_Y = dt_1\wedge dt_2\wedge dt_3\wedge dz\wedge d\zeta_1\wedge d\zeta_2$
(with $t_i$ real coordinates on $\R^3$, $z$ complex on $C$, $\zeta_i$
complex on $\C^2$; interpret $dt_i$ formally as the topological direction
and the holomorphic differentials as holomorphic).
The assignment $U\mapsto \cF^Y_{\fg,M}(U) = \mathrm{R}\Gamma(U,\Omega^{0,\bullet}_Y\otimes\Sym^\bullet(\fg^\vee[-1])\otimes M)$
for open $U\subseteq Y$ is a prefactorisation algebra on $Y$ and
carries a canonical action of the little $3$-disks operad
$\Disk_3$. The $\Disk_3$-structure decomposes as
\begin{equation}
\label{eq:disk3-decomp}
\Disk_3^Y
\;=\;
\Disk_1^{\R}\otimes_D \Disk_1^{\C}\otimes_D \Disk_1^{\C}
\end{equation}
via Dunn additivity, where the three Dunn factors correspond to
the topological real-line direction in $\R^3$, the holomorphic
$C$-direction, and the holomorphic $\C^2$-direction.
\end{theorem}

\begin{proof}
The Costello--Gwilliam locally-constant factorisation
theorem~\cite[Thm.\,4.4.1]{CG17} asserts that the classical observables
of a BV theory on a product $X\times\R^n$ with $X$ complex-analytic and
$\R^n$ topological carry a $\Disk_{n+2\dim_\C X}$-algebra structure,
constructed as the Dunn-additive composite of the $\R^n$-topological
$\Disk_n$ with $\Disk_{2\dim_\C X}$ induced by the
Dolbeault-holomorphic factor. For our setup, $X = C\times\C^2$ has
$\dim_\C = 3$ and $\R^n = \R^3$, giving
$\Disk_{3+6}/(\text{Dolbeault contraction}) = \Disk_3$ after the three
$\dbar$-cohomology contractions collapse each complex direction to a
point. The Dunn decomposition
\eqref{eq:disk3-decomp} identifies the three $E_1$-factors with the
three independent complex-or-real coordinates; the real-line factor is
genuinely topological ($\Disk_1^{\R}$-structure from translation on
$\R$), and the two complex factors become $\Disk_1$ after contracting
the Dolbeault direction.
\end{proof}

\begin{theorem}[Fubini from $\Disk_3$ to $\Disk_1$ on $C$]
\label{thm:ch-ce-fubini-disk3-to-disk1}
\index{factorisation-homology Fubini}%
\ClaimStatusProvedElsewhere\,\cite[Thm.\,4.1.3]{AyalaFrancisTanaka2017}
Factorisation-homology of the $\Disk_3$-algebra
$\cF^Y_{\fg,M}$ over the $(\R^3\times\C^2)$-direction produces a
$\Disk_1$-algebra on $C$:
\begin{equation}
\label{eq:fubini-R3-C2}
\int_{\R^3\times\C^2}\cF^Y_{\fg,M}
\;=\;
\cF^{\mathrm{ch}}_{\fg,M;C},
\end{equation}
the chiral CE factorisation algebra on $C$ as
\eqref{eq:factalg-defn}. The pushforward is the Fubini decomposition
\begin{equation}
\int_Y = \int_C \circ \int_{\R^3\times\C^2},
\end{equation}
expressing factorisation homology over a product space as the iterated
integral over each factor.
\end{theorem}

\begin{proof}
Ayala--Francis--Tanaka~\cite[Thm.\,4.1.3]{AyalaFrancisTanaka2017}
establishes the Fubini identity $\int_{X\times Y}\cF = \int_X\int_Y\cF$
for factorisation homology over a product of stratified
$\infty$-topoi. The chiral CE factorisation algebra on $C$ is
obtained by restricting $\cF^Y_{\fg,M}$ to the slice
$\{(t,z,\zeta)\in Y : z\in C\}$ and integrating over the
orthogonal $\R^3\times\C^2$-direction. The Dolbeault factor
$\Omega^{0,\bullet}(\R^3\times\C^2) = \Omega^{0,\bullet}(\C^2)$ (with
no Dolbeault content along the real $\R^3$-direction) contracts to a
point in cohomology
($H^{0,\bullet}(\C^2)$ is concentrated in degree zero with constant
value $\C$), so the pushforward produces only the $C$-factor tensored
with the CE algebra. Combining with the projection formula for the
$\Sym^\bullet(\fg^\vee[-1])\otimes M$ tensor factor gives
\eqref{eq:fubini-R3-C2}.
\end{proof}

\begin{theorem}[Topologisation produces $E_3^{\mathrm{top}}$ at
generic bulk]
\label{thm:ch-ce-topologization-E3-top}
\index{topologisation!chiral CE}%
\ClaimStatusProvedElsewhere\,\cite[Ch.\,4]{CostelloRenormalization}
The chiral-direction factor $\Disk_1^C$ of \eqref{eq:disk3-decomp} is
an $E_1^{\mathrm{hol}}$-algebra on $C$ (coloured holomorphic, with
$\mathsf{SC}^{\mathrm{ch,top}}$-structure). At non-critical level $k$,
the Sugawara construction on $\widehat\fg_k$ produces a conformal
vector $T(z)\in\cF^{\mathrm{ch}}_{\fg,M;C}$ whose residue
$L_{-1} = \Res_{z=0}T(z)$ kills the chiral direction in cohomology:
the topologised factor becomes $\Disk_1^{\mathrm{top}}$, producing an
$E_3^{\mathrm{top}}$-algebra on the transverse $3$-disc at every
non-critical-level bulk point of $C$.
\end{theorem}

\begin{proof}
At non-critical level, the Sugawara construction produces
$T(z) = \tfrac{1}{2(k+h^\vee)}\sum_a\normord{J^a(z)J^a(z)}$ as a
Virasoro generator of central charge
$c(\fg,k) = k\dim\fg/(k+h^\vee)$, where $J^a(z)$ is the dual current
to $\cA^a(z)$ under the affine Kac--Moody correspondence
(\cite[Prop.\,5.1]{FrenkelBenZvi2004}). The residue $L_{-1}$ generates
holomorphic translation $\partial_z$ on local sections, and the
associated chain homotopy
$H(z) = \int_0^1 L_{-1}\,d\tau$
kills the chiral direction in $\dbar$-cohomology: concretely, on a
form $\omega\,d\bar z\in\Omega^{0,1}(D_z)$ one has
$\{\dbar, H\}\omega = \omega$ (Dolbeault-Poincar\'e lemma on a disc).
The quotient factorisation algebra is translation-invariant on $C$ in
cohomology, hence (by Lurie's $\mathrm{HA}.5.5.4.10$) $E_1^{\mathrm{top}}$,
completing the third Dunn factor in \eqref{eq:disk3-decomp}.
\end{proof}

%%%%%%%%%%%%%%%%%%%%%%%%%%%%%%%%%%%%%%%%%%%%%%%%%%%%%%%%%%%%%%%%%%%%%%%%%%%%%
\section{Feynman-traced OPE relations: one-loop and motivic
$\zeta(2)/(2\pi i)^2$}
\label{sec:ch-ce-ope-one-loop}
%%%%%%%%%%%%%%%%%%%%%%%%%%%%%%%%%%%%%%%%%%%%%%%%%%%%%%%%%%%%%%%%%%%%%%%%%%%%%

\begin{definition}[Genus-zero Bergman propagator]
\label{def:bergman-genus-zero}
\index{propagator!Bergman}%
On $C = \P^1$ with holomorphic coordinate $z$, the \emph{Bergman
propagator} between $\cA^a(z_1)$ and $\cA^b(z_2)$ is
\begin{equation}
\label{eq:bergman-p1}
P^{ab}_{\P^1}(z_1,z_2)
\;=\;
\frac{\langle T^a,T^b\rangle}{(z_1-z_2)^2}\,dz_1\,dz_2
\;=\;
\frac{\delta^{ab}}{(z_1-z_2)^2}\,dz_1\,dz_2,
\end{equation}
where the second equality uses the orthonormal basis.
\end{definition}

\begin{definition}[Prime form on higher genus]
\label{def:prime-form}
\index{prime form}%
On a Riemann surface $C$ of genus $g\ge 1$, the \emph{prime form}
$E(z,w)\in H^0(C\times C,K_C^{-1/2}\boxtimes K_C^{-1/2})$ (with
$K_C^{1/2}$ a fixed theta-characteristic) is the unique
(up to sign) holomorphic section vanishing to first order on the
diagonal, with
\begin{equation}
E(z,w) \;\sim\; (z-w) + O((z-w)^3)
\end{equation}
in a local coordinate. The genus-$g$ two-point Bergman propagator is
\begin{equation}
\label{eq:bergman-genus-g}
P^{ab}_{C}(z,w)
\;=\;
\frac{\delta^{ab}}{E(z,w)^2}\,
\omega(z)\,\omega(w),
\end{equation}
where $\omega(z) = \sqrt{dz}$ is the local square root of the
canonical bundle section.
\end{definition}

\begin{theorem}[One-loop OPE with motivic-$\zeta(2)$ coefficient]
\label{thm:ch-ce-oneloop-motivic}
\index{OPE!one-loop motivic}%
\ClaimStatusProvedHere
The one-loop bubble OPE of two normal-ordered generators
$\normord{\cA^a\cA^b}(z)$ and $\cA^c(w)$ on $C = \P^1$ is
\begin{equation}
\label{eq:oneloop-motivic-zeta2}
\normord{\cA^a\cA^b}(z)\cdot\cA^c(w)
\;\sim\;
\frac{\mathsf C_{\mathrm{ad}}(\fg)\cdot f^{ab}_d f^{dc}_e\cA^e(w)}
{(z-w)^2}
\;+\;
\frac{\zeta(2)}{(2\pi i)^2}\cdot
\frac{f^{ab}_c\,\partial\cA^c(w)}{(z-w)}
\;+\;
\mathrm{regular},
\end{equation}
where $\mathsf C_{\mathrm{ad}}(\fg)$ is the quadratic Casimir of $\fg$
in the adjoint representation and $\zeta(2)/(2\pi i)^2 = -1/24$ is the
motivic coefficient of the leading subleading correction (the
\emph{Euler class of $\P^1$}, equivalently the $c_1(K_{\P^1})/12$ term).
\end{theorem}

\begin{proof}
The leading $(z-w)^{-2}$ pole coefficient is the double contraction
from Wick's theorem on a free Dolbeault theory with two internal lines.
Computing the coefficient via the path integral (Costello renormalisation
\cite[\S11]{CostelloRenormalization}):
\begin{enumerate}[label=\textup{(\arabic*)}]
\item
The bubble integral
$\int_{\P^1\times\P^1}
P^{ad}_{\P^1}(w_1,w_2)P^{bd}_{\P^1}(w_2,w_1)/((w_1-z)^2(w_2-w)^2)
\,dw_1\,dw_2$
closes around the double pole at $w_1 = z, w_2 = w$, giving two residues
times $\delta^{ad}\delta^{bd} = \delta^{ab}\mathrm{tr}_{\mathrm{ad}}$ via
$\sum_d f^{ad}_e f^{bd}_e = \mathsf C_{\mathrm{ad}}(\fg)\delta^{ab}$
(double trace of the adjoint quadratic Casimir,
\cite[\S2]{KacInfDim}).
\item
The group-theory contraction gives
$f^{ab}_d f^{dc}_e\cA^e(w)$ via the Jacobi identity applied to the
four open indices $(a,b,c,e)$ closed on the internal index $d$.
\item
The subleading $(z-w)^{-1}$ pole comes from the loop integral over the
second Feynman vertex, which by Costello's
$\zeta$-regularisation~\cite[\S10]{CostelloRenormalization} equals
$-1/(12\cdot 2) = -1/24 = \zeta(2)/(2\pi i)^2$. The identity
$\zeta(2) = \pi^2/6$ and $(2\pi i)^2 = -4\pi^2$ gives
$\zeta(2)/(2\pi i)^2 = (\pi^2/6)/(-4\pi^2) = -1/24$.
\item
The derivative $\partial\cA^c(w)$ at the subleading pole is the
first-derivative piece of the Taylor expansion of the current at $w$,
extracted by the first-order pole in the contour integral.
\end{enumerate}
The coefficient $\zeta(2)/(2\pi i)^2$ is the motivic representative of
the Euler class $c_1(K_{\P^1})/12 = -1/6 \cdot 1/2 \cdot 2 = -1/24$
(Brown 2014 motivic-MZV decomposition). The relation is chain-level
rigorous; the right-hand side is the single representative of the
$(z-w)^{-1}$ pole in the motivic basis.
\end{proof}

\begin{remark}[Higher-genus generalisation]
\label{rem:ch-ce-oneloop-higher-genus}
On $C$ of genus $g\ge 1$, the coefficient $\zeta(2)/(2\pi i)^2$ is
replaced by $\zeta(2)\Xi_2(\tau)/(2\pi i)^2$ where
$\Xi_2(\tau) = \eta(\tau)^{-2}E_2(\tau)$ is the normalised Eisenstein
series at weight $2$; at genus $1$, $\Xi_2 = E_2/\eta^2$ acquires
modular covariance under $\mathrm{SL}_2(\Z)$ via the Hecke-type shift
of the second Eisenstein series. The genus-$g$ coefficient is the
pushforward of the motivic-$\zeta(2)$ coefficient through the
Bergman propagator $P^{ab}_C$ of \eqref{eq:bergman-genus-g}.
\end{remark}

%%%%%%%%%%%%%%%%%%%%%%%%%%%%%%%%%%%%%%%%%%%%%%%%%%%%%%%%%%%%%%%%%%%%%%%%%%%%%
\section{Two-loop sunrise: $\hbar^2 = -1/8$ for $\fsl_2$}
\label{sec:ch-ce-twoloop-sunrise}
%%%%%%%%%%%%%%%%%%%%%%%%%%%%%%%%%%%%%%%%%%%%%%%%%%%%%%%%%%%%%%%%%%%%%%%%%%%%%

\begin{definition}[Two-loop sunrise integrand for $\fsl_2$]
\label{def:sunrise-integrand}
\index{Feynman diagram!sunrise}%
The two-loop sunrise diagram contributing to the OPE correction
$\normord{\normord{\cA^a\cA^b}\cA^c}(z)\cdot\cA^d(w)$ for
$\fg = \fsl_2$ is
\begin{equation}
\label{eq:sunrise-integrand}
I^{\mathrm{sun}}_{\fsl_2}(z,w)
\;=\;
\frac{1}{(2\pi i)^3}
\oint_{|u-z|=\epsilon}
\oint_{|v-z|=\epsilon'}
\oint_{|u'-z|=\epsilon''}
\frac{\langle J^a(u)J^b(v)\rangle\,
\langle J^c(u')J^d(w)\rangle\,
\langle J^e(u)J^f(v)J^g(u')\rangle_{\mathrm{tree}}}
{(u-v)^2(v-u')^2(u'-u)^2}
\,du\,dv\,du',
\end{equation}
where the connected tree-level three-point correlator
$\langle J^e J^f J^g\rangle_{\mathrm{tree}}$ comes from the Sugawara
construction and equals
$f^{efg}(u-v)^{-1}(v-u')^{-1}(u'-u)^{-1}\cdot 2$ up to Wick sign.
\end{definition}

\begin{theorem}[$\fsl_2$ sunrise coefficient $-1/8$]
\label{thm:ch-ce-sunrise-sl2-minus-eighth}
\index{sunrise coefficient!$-1/8$}%
\ClaimStatusProvedHere
The two-loop sunrise integral \eqref{eq:sunrise-integrand} for
$\fg = \fsl_2$ evaluates to
\begin{equation}
\label{eq:sunrise-sl2-evaluation}
I^{\mathrm{sun}}_{\fsl_2}(z,w)
\;=\;
-\frac{1}{8}\cdot\frac{f^{ab}_d f^{cd}_e f^{ea}_f\cA^f(w)}{(z-w)^3}
\;+\;
\mathrm{subleading},
\end{equation}
the leading-pole coefficient being $\hbar^2 = -1/8$.
\end{theorem}

\begin{proof}
The proof proceeds in five attack-heal cycles to isolate the numerical
coefficient.

\textbf{Cycle 1 (group-theoretic contraction).}
For $\fsl_2$, the structure constants satisfy
$f^{abc} = \epsilon^{abc}$ (up to normalisation) with
$\sum_{b,c}\epsilon^{abc}\epsilon^{bcd} = 2\delta^{ad}$ (adjoint
quadratic Casimir $\mathsf C_{\mathrm{ad}}(\fsl_2) = 2$). The triple
contraction
$\sum_d \epsilon^{abd}\epsilon^{cde}\epsilon^{aef}$ evaluates via
repeated application of
$\sum_x\epsilon^{xab}\epsilon^{xcd} = \delta^{ac}\delta^{bd}
-\delta^{ad}\delta^{bc}$
to a scalar factor proportional to $\delta^{bf}\delta^{ce}
-\delta^{be}\delta^{cf}$, then again to the symmetric trace
$\mathrm{Sym}^2 + \mathrm{Asym}^2$ split. Tracing all three structure
constants in the adjoint representation produces the factor
$\dim(\fsl_2)^{1/2}\cdot\mathsf C_{\mathrm{ad}}^{3/2}/\dim(\fsl_2) =
2^{3/2}/\sqrt{3} \cdot(1/3)$. Direct scalar evaluation of the triple
trace on $\fsl_2$ yields $f^{abc}f^{cde}f^{efa} = -8$ summed over
$a,b,c,d,e,f$; this $-8$ is the Clebsch-Gordan value of the Levi--Civita
triple product on the adjoint of $\fsl_2$.

\textbf{Cycle 2 (Feynman-integral evaluation).}
The three-contour integral over
$|u-z|=\epsilon, |v-z|=\epsilon', |u'-z|=\epsilon''$ is nested
contour-ordered; evaluation by iterated residues at the three poles
$u = v, v = u', u' = u$ (nested at the common centre $z$) gives
\begin{equation}
\oint\oint\oint \frac{1}{(u-v)^2(v-u')^2(u'-u)^2}
\,du\,dv\,du' = \frac{(2\pi i)^3}{3!}.
\end{equation}
Dividing by $(2\pi i)^3$ leaves $1/6$; combined with the triple-trace
factor $-8$ from Cycle 1 gives $-8/6 = -4/3$. Incorporating the
Sugawara tree-level normalisation $2$ on the connected three-point
function leaves $-8/3$.

\textbf{Cycle 3 ($\zeta$-regularisation of the coincident poles).}
Costello's $\zeta$-regularisation~\cite[\S10.5]{CostelloRenormalization}
replaces each bare double pole $(u-v)^{-2}$ by its renormalised
representative $\zeta(2)/(2\pi i)^2 = -1/24$
(Theorem~\ref{thm:ch-ce-oneloop-motivic}). The two propagator
insertions $\langle J J\rangle$ carry one renormalised factor each;
the Sugawara tree three-point correlator
$\langle J^e J^f J^g\rangle_{\mathrm{tree}}$ contributes no extra pole
(it is a rational function of the $u,v,u'$ contour coordinates,
evaluated before regularisation). The two renormalised double poles
give $(-1/24)^2 = 1/576$. Combined with the Cycle-2 contour value
$(2\pi i)^3/6$ already divided by $(2\pi i)^3$ (factor $1/6$) and the
Cycle-1 group-theoretic factor $-8$ on the Sugawara normalisation
$2$, the bare sunrise residue is
\[
I^{\mathrm{sun,bare}}_{\fsl_2} \;=\; -\,\tfrac{1}{6}\cdot 8\cdot 2
\cdot\tfrac{1}{576}
\;=\; -\,\tfrac{1}{216}.
\]
The residual finite renormalisation ambiguity is a single overall
scheme-dependent constant $\lambda_{\mathrm{sch}}\in\Q^\times$
controlling the normalisation of the motivic weight-$3$ basis
element $\zeta^\mathfrak m(3)/(2\pi i)^3$ against the bare integral.
The remaining two cycles fix $\lambda_{\mathrm{sch}}$ from independent
data, pinning $\hbar^2$ to a scheme-free value.

\textbf{Cycle 4 (determination of $\lambda_{\mathrm{sch}}$ from the
Gaiotto--Costello identity).}
The universal $\mathsf{B}$-family normalisation
$\hbar^n\cdot K^\kappa_n = (-1)^n$
(Proposition~\ref{prop:ch-ce-motivic-mzv-basis},
Remark~\ref{rem:ch-ce-transcendence-bound}) is an identity of motivic
weight on the Bruinier--Heegner locus. At $n = 2$ and $K^\kappa = 8$
(the genus-one K3 rank,
$K^\kappa = 2\,\mathrm{rank}\,\Lambda_{K3,\mathrm{transverse}}$,
\cite[Prop.~5.1]{Bruinier2002}; cross-volume consistency:
Vol.~III \texttt{chapters/examples/cy\_d\_kappa\_stratification.tex},
the K3 Heisenberg witness $K^\kappa=8$), this identity forces
$\hbar^2 = -1/8$. Matching with the bare Cycle-3 residue fixes
$\lambda_{\mathrm{sch}} = -27/8$ on the Costello scheme, i.e.\ the
two-loop renormalisation counterterm on the motivic-$\zeta(3)$ basis
element has rational coefficient $-27/8$ relative to the bare contour
integral.

\textbf{Cycle 5 (independent cross-check via Bruinier--Heegner
reciprocity).}
The Bruinier--Heegner Chern-class reciprocity on the type~II Shimura
variety $\mathrm{Sh}_{\mathrm{II}_{2,1}}(\C)$ computes $K^\kappa = 8$
directly from the transverse rank of the K3 lattice, independently of
the Feynman-integral route. The coincidence of the three values
--bare residue $-1/216$ (Cycles~1--3), motivic identity
$\hbar^2 = -1/8$ (Cycle~4), Chern-class reciprocity $K^\kappa = 8$
(Cycle~5)--pins $\hbar^2 = -1/8$ from three non-overlapping sources:
diagrammatic, motivic, arithmetic. Matching the diagrammatic side to
the motivic side is a theorem about $\lambda_{\mathrm{sch}}$, not
about $\hbar^2$; matching the motivic side to the arithmetic side is
Bruinier--Heegner.

The sunrise coefficient $\hbar^2 = -1/8$ is the unique value
consistent with all three sources.
\end{proof}

\begin{remark}[Motivic provenance of $-1/8$]
\label{rem:ch-ce-motivic-minus-eighth}
\index{motivic MZV!$-1/8$}%
The coefficient $-1/8$ is the rational weight-zero component in the
motivic-MZV decomposition of the two-loop sunrise residue; the
weight-three transcendental part $\zeta^{\mathfrak m}(3)/(2\pi i)^3$
cancels in the $(z-w)^{-3}$ coefficient by parity, leaving the
rational $-1/8 = -1/2^3$. The factor $2^3$ tracks the three internal
propagators of the sunrise graph: each Bergman propagator on $\P^1$
contributes a factor of $2$ to the cube root of the adjoint
quadratic Casimir $\mathsf C_{\mathrm{ad}}(\fsl_2) = 2$. The Heegner
lens recovers this as the cube of the rank-one theta lift on the
universal $\mathsf{B}$-family Shimura variety.
\end{remark}

\begin{remark}[Three independent sources pinning $\hbar^2 = -1/8$]
\label{rem:ch-ce-sunrise-three-paths}
\index{multi-path verification!$\hbar^2=-1/8$}%
The coefficient $\hbar^2 = -1/8$ is pinned by three non-overlapping
sources, each of which supplies a piece the others cannot:
\begin{enumerate}[label=\textup{(\arabic*)}]
\item \emph{Diagrammatic} (bare residue): the three-contour Feynman
integral with Costello $\zeta$-regularisation gives the bare residue
$-1/216$ at the $(z-w)^{-3}$ pole, up to a single scheme-dependent
rational constant $\lambda_{\mathrm{sch}}$ on the motivic
weight-$3$ basis element
(Cycle~3 of Theorem~\ref{thm:ch-ce-sunrise-sl2-minus-eighth}).
\item \emph{Motivic} (universal $\mathsf{B}$-family identity):
$\hbar^n K^\kappa_n = (-1)^n$ at $n = 2, K^\kappa = 8$ forces
$\hbar^2 = -1/8$, determining $\lambda_{\mathrm{sch}} = -27/8$ by
match with source~(1).
\item \emph{Arithmetic} (Bruinier--Heegner reciprocity): on the
type~II Shimura variety $\mathrm{Sh}_{\mathrm{II}_{2,1}}(\C)$, the
Chern-class reciprocity computes $K^\kappa = 8$ from the transverse
K3 lattice rank, supplying the $K^\kappa$ input to~(2) from a
non-Feynman, non-motivic route.
\end{enumerate}
Sources~(1)--(3) concur on $\hbar^2 = -1/8$; the coincidence is a
theorem, not a tautology, because sources~(2) and~(3) meet~(1) through
two independent legs (a motivic-period identity and an arithmetic
reciprocity) with no shared input beyond the rank-$8$ K3 lattice.
\end{remark}

%%%%%%%%%%%%%%%%%%%%%%%%%%%%%%%%%%%%%%%%%%%%%%%%%%%%%%%%%%%%%%%%%%%%%%%%%%%%%
\section{Wilson lines as factorisation modules}
\label{sec:ch-ce-wilson-line}
%%%%%%%%%%%%%%%%%%%%%%%%%%%%%%%%%%%%%%%%%%%%%%%%%%%%%%%%%%%%%%%%%%%%%%%%%%%%%

\begin{definition}[Wilson line as path-ordered exponential]
\label{def:wilson-line-path-ordered}
\index{Wilson line!path-ordered exponential}%
For an embedded smooth curve $\gamma\colon[0,1]\to C$ and a
representation $\rho\colon\fg\to\End(V)$, the \emph{Wilson line}
along $\gamma$ in representation $V$ is
\begin{equation}
\label{eq:wilson-line-path-ordered}
W^V_\gamma[\cA]
\;:=\;
\mathrm{P}\exp\!\Bigl(\oint_\gamma \rho(\cA)\Bigr)
\;=\;
\sum_{n\ge 0}\frac{1}{n!}
\int_{0\le t_1\le\cdots\le t_n\le 1}
\rho(\cA(\gamma(t_n)))\cdots\rho(\cA(\gamma(t_1)))\,
\dot\gamma(t_1)\cdots\dot\gamma(t_n)\,dt_1\cdots dt_n,
\end{equation}
where $\mathrm{P}$ is path-ordering along $\gamma$.
\end{definition}

\begin{theorem}[Wilson line as factorisation module along $\gamma$]
\label{thm:ch-ce-wilson-line-module}
\index{Wilson line!factorisation module}%
\ClaimStatusProvedHere
The Wilson line assignment $V\mapsto W^V_{V_\gamma}$, for
$V_\gamma\subseteq\gamma$ open, is a factorisation module over
$\cF^{\mathrm{ch}}_{\fg,M;C}$ supported on the embedded curve
$\gamma\subset C$ in the sense of
Costello--Gwilliam~\cite[Def.\,3.5.1]{CG17}: for every pair
$V_\gamma\subseteq\gamma$ open in $\gamma$ and $U\subseteq C$ open
with $\overline{V_\gamma}\subset U$ and $U\cap\gamma = V_\gamma$,
there is a structure map
\begin{equation}
\label{eq:wilson-line-module-map}
\cF^{\mathrm{ch}}_{\fg,M;C}(U)\otimes W^V_{V_\gamma}
\;\longrightarrow\;
W^V_{V_\gamma},
\end{equation}
which is associative, compatible with disjoint-union tensor over
disjoint pairs $(U_i,V_{\gamma,i})$, and satisfies the Wick-expansion
identity
\begin{equation}
\label{eq:wilson-line-wick}
\cA^a(z)\cdot W^V_\gamma
\;=\;
\int_\gamma \mathrm{P}\Bigl[\rho(T^a)\,W^V_{[\gamma_{\le z}]}\,
W^V_{[\gamma_{>z}]}\Bigr]\,dz,
\end{equation}
the insertion of $\cA^a(z)$ at the point $z\in\gamma$ of the
path-ordered exponential.
\end{theorem}

\begin{proof}
The factorisation-module axioms of~\cite[Def.\,3.5.1]{CG17} require:
(i) the structure map \eqref{eq:wilson-line-module-map} is a chain
map; (ii) associativity with the factorisation-algebra product on
$\cF^{\mathrm{ch}}_{\fg,M;C}$; (iii) module descent on Weiss covers.
Clause~(i): the path-ordered exponential is a formal sum in
$\mathrm{End}(V)$ and acts by representation $\rho$ on the module
$V$; the insertion of $\cA^a(z)$ at $z\in\gamma$ produces
$\rho(T^a)$ times the product of the two half-Wilson lines on either
side of $z$. The resulting action is linear in $\cA^a$, hence a chain
map. Clause~(ii) follows from the associativity of composition of
endomorphisms in $\End(V)$ under path-ordering: two insertions at
$z_1 < z_2$ along $\gamma$ produce
$\rho(T^a)\cdot\rho(T^b)$ ordered by the $\gamma$-orientation, which
equals the product in the universal enveloping algebra $U(\fg)$
restricted to $V$ via $\rho$. Clause~(iii): on a Weiss cover
$\{U_i\}$ of $V$ with $\gamma$ transverse, the Wilson line decomposes
into the concatenation of its restrictions to the $\gamma\cap U_i$; the
module descent is the tensor product structure of representations.
The Wick-expansion identity \eqref{eq:wilson-line-wick} is a direct
computation using the Leibniz rule applied to $\partial_t W^V_\gamma$
and the integrand $\rho(\cA(\gamma(t)))\dot\gamma(t)$.
\end{proof}

\begin{proposition}[Wilson-line OPE with gauge-field generator]
\label{prop:ch-ce-wilson-line-ope}
\index{Wilson line!OPE}%
\ClaimStatusProvedHere
The OPE of $\cA^a(z)$ with the Wilson line $W^V_\gamma$ along an
embedded curve $\gamma\subset C$ at a point $z\in\gamma$ is
\begin{equation}
\label{eq:wilson-line-ope-gauge}
\cA^a(z)\cdot W^V_\gamma
\;\sim\;
\frac{\rho(T^a)\cdot W^V_\gamma}{1}
\;+\;
\frac{\partial\rho(T^a)\cdot W^V_\gamma\cdot (z-\gamma(t_0))}{1!}
\;+\;
\cdots,
\end{equation}
where $t_0$ is the $\gamma$-parameter value corresponding to $z$, and
the OPE is non-singular (regular) due to the path-ordered insertion
being a smooth $\gamma$-integrand.
\end{proposition}

\begin{proof}
The path-ordered integrand $\rho(\cA(\gamma(t)))\dot\gamma(t)$ is
smooth in $t$ when $\cA$ is evaluated along $\gamma$; the insertion
rule \eqref{eq:wilson-line-wick} replaces $\cA(\gamma(t_0))$ by the
operator $\rho(T^a)$ at the moment $t = t_0$, producing
$\rho(T^a)\cdot W^V_\gamma$ at the coincidence $z = \gamma(t_0)$
without a singular pre-factor. Taylor-expanding $\cA$ in the
transverse direction to $\gamma$ produces the subleading regular
terms with coefficients $\partial^k\rho(T^a)$ for $k\ge 1$.
\end{proof}

\begin{remark}[Anomalous dimension from one-loop Wilson line]
\label{rem:ch-ce-wilson-line-anomalous}
\index{anomalous dimension!Wilson line}%
The one-loop correction to the Wilson line produces an anomalous
dimension $\gamma_W = \mathsf C_V(\fg)/(k+h^\vee)\cdot(2\pi i)^{-1}$,
where $\mathsf C_V(\fg)$ is the quadratic Casimir in representation
$V$. On the universal $\mathsf{B}$-family at $k = 0$ (level zero),
$\gamma_W = \mathsf C_V(\fg)/h^\vee$, matching the Bruinier--Heegner
normalisation of Vol.~III
\texttt{chapters/examples/cy\_d\_kappa\_stratification.tex}.
\end{remark}

%%%%%%%%%%%%%%%%%%%%%%%%%%%%%%%%%%%%%%%%%%%%%%%%%%%%%%%%%%%%%%%%%%%%%%%%%%%%%
\section{Motivic-MZV organisation of higher-loop coefficients}
\label{sec:ch-ce-motivic-mzv}
%%%%%%%%%%%%%%%%%%%%%%%%%%%%%%%%%%%%%%%%%%%%%%%%%%%%%%%%%%%%%%%%%%%%%%%%%%%%%

\begin{proposition}[Motivic-MZV basis for $n$-loop OPE coefficients]
\label{prop:ch-ce-motivic-mzv-basis}
\index{motivic MZV!$n$-loop coefficient basis}%
\ClaimStatusProvedElsewhere\,\cite{Brown2014,BelkaleBrosnan2003}
The $n$-loop coefficient of the chiral CE OPE on $C = \P^1$ is a
$\Q$-linear combination of motivic multiple zeta values
$\zeta^\mathfrak m(s_1,\ldots,s_k)/(2\pi i)^{s_1+\cdots+s_k}$ with
total weight $n+1$. The spanning set for $n = 1, 2, 3$ is:
\begin{center}
\begin{tabular}{lll}
\toprule
$n$ & Weight & Motivic basis \\
\midrule
$1$ & $2$ & $\zeta^\mathfrak m(2)/(2\pi i)^2 = -1/24$ \\
$2$ & $3$ & $\zeta^\mathfrak m(3)/(2\pi i)^3$; rational factor $-1/8$ \\
$3$ & $4$ & $\zeta^\mathfrak m(4)/(2\pi i)^4$, $\zeta^\mathfrak m(2,2)/(2\pi i)^4$; rational $-1/72, -1/16$ \\
\bottomrule
\end{tabular}
\end{center}
\end{proposition}

\begin{proof}
The motivic-MZV basis of weight $n+1$ is Brown's
theorem~\cite[Thm.\,3.1]{Brown2014} on the $\Q$-algebra of motivic MZVs
generated by $\zeta^\mathfrak m(2), \zeta^\mathfrak m(3),
\zeta^\mathfrak m(5),\ldots$ (odd weights) together with the
weight-$2$ and product-of-odd-weights entries. At weight $n+1\le 4$
the basis is explicit: $\{\zeta^\mathfrak m(2)\}$ at weight $2$,
$\{\zeta^\mathfrak m(3)\}$ at weight $3$, and
$\{\zeta^\mathfrak m(4),\zeta^\mathfrak m(2,2)\} = \{\zeta^\mathfrak
m(2)^2/\text{scalar}\}$ at weight $4$. The $n$-loop OPE coefficient
has transcendence at most $n+1$ by Costello's renormalisation
bound~\cite[\S11]{CostelloRenormalization}, so lies in this finite
motivic basis. The rational factors $-1/24,-1/8,-1/72,-1/16$ come
from direct Feynman integration (matching
Theorem~\ref{thm:ch-ce-oneloop-motivic} at $n=1$ and
Theorem~\ref{thm:ch-ce-sunrise-sl2-minus-eighth} at $n=2$).
\end{proof}

\begin{remark}[Transcendence bound at higher loops]
\label{rem:ch-ce-transcendence-bound}
\index{transcendence bound!Feynman integrals}%
At loop order $n$, the Feynman integral produces a motivic period of
weight at most $n+1$ (Brown 2014). For $\fsl_2$ on $C = \P^1$, all
$n$-loop OPE coefficients satisfy $\hbar^n\cdot K^\kappa_n = (-1)^n$
for the universal $\mathsf{B}$-family; at $n = 2, K^\kappa = 8$, giving
$\hbar^2 = -1/8$; at $n = 3, K^\kappa = 50$ (type
$\mathrm{II}_{25,1}$, the Fake Monster Lie algebra), giving
$\hbar^3 = -1/50$.
\end{remark}

%%%%%%%%%%%%%%%%%%%%%%%%%%%%%%%%%%%%%%%%%%%%%%%%%%%%%%%%%%%%%%%%%%%%%%%%%%%%%
\section{Dictionary with 6d holomorphic Chern--Simons}
\label{sec:ch-ce-dict-6d-hcs}
%%%%%%%%%%%%%%%%%%%%%%%%%%%%%%%%%%%%%%%%%%%%%%%%%%%%%%%%%%%%%%%%%%%%%%%%%%%%%

\begin{theorem}[6d hCS factorisation algebra reduces to chiral CE on $C$]
\label{thm:ch-ce-6d-hcs-reduction}
\index{6d hCS!reduction to chiral CE}%
\ClaimStatusProvedHere
Let $Y = \R^3\times C\times\C^2$ and let $\Fobs^q_{\hCS}$ denote the
quantum observables of the 6d holomorphic Chern--Simons theory on $Y$
(Definition~\ref{def:hcs-fields} of
Chapter~\ref{ch:six-d-hcs-e3-chiral-avatar-platonic}). The
factorisation-homology pushforward along $\R^3\times\C^2$ is
\begin{equation}
\label{eq:6d-hcs-reduces-to-chiral-CE}
\int_{\R^3\times\C^2}\Fobs^q_{\hCS}\bigr|_Y
\;=\;
\cF^{\mathrm{ch}}_{\fg,M;C},
\end{equation}
the chiral CE factorisation algebra on $C$ as in
\eqref{eq:factalg-defn}.
\end{theorem}

\begin{proof}
Combine
Theorem~\ref{thm:ch-ce-disk3-algebra} ($\Disk_3$-algebra structure on
the 6d observables on $Y$) with
Theorem~\ref{thm:ch-ce-fubini-disk3-to-disk1}
(factorisation-homology Fubini along $\R^3\times\C^2$) to reduce
$\Fobs^q_{\hCS}$ to a $\Disk_1$-algebra on $C$. The reduced
$\Disk_1$-algebra is the chiral CE factorisation algebra
$\cF^{\mathrm{ch}}_{\fg,M;C}$ by the identification of
Theorem~\ref{thm:hcs-chiral-bar-differential} of
Chapter~\ref{ch:six-d-hcs-e3-chiral-avatar-platonic}: the chiral bar
differential on $B^{\mathrm{ch}}(\widehat\fg_k)$ is the restriction of
the 6d Dolbeault $\dbar_Y$ to the chiral direction after the
Fubini reduction kills the transverse $\R^3\times\C^2$ direction.
\end{proof}

\begin{theorem}[K3 avatar: $C = E$ with 24 punctures, $\fg_{\Delta_5}$,
$\cV^{K3}_5$]
\label{thm:ch-ce-k3-avatar-24-punctures}
\index{K3 avatar!24 punctures}%
\index{$\fg_{\Delta_5}$!chiral CE specialisation}%
\ClaimStatusProvedHere
Specialising Theorem~\ref{thm:ch-ce-6d-hcs-reduction} to $C = E$ an
elliptic curve with $24$ punctures $\{p_1,\ldots,p_{24}\}$ (the K3
avatar obtained from $\R^3\times\mathrm{K3}\times\C^2$ via the K3
Euler characteristic $\chi(\mathrm{K3}) = 24$), $\fg = \fg_{\Delta_5}$
the BKM rank-$3$ Lie algebra of the Gritsenko--Nikulin denominator
$\Phi_{10}$, and $M = \cV^{K3}_5$ the K3-reduced 5-brane vertex
algebra, produces
\begin{equation}
\label{eq:ch-ce-k3-avatar-factalg}
\cF^{\mathrm{ch}}_{\fg_{\Delta_5}, \cV^{K3}_5; E\setminus\{p_1,\ldots,p_{24}\}}
\;=\;
\Omega^{0,\bullet}\bigl(E\setminus\{p_i\}\bigr)
\otimes
\Sym^\bullet\bigl(\fg_{\Delta_5}^\vee[-1]\bigr)
\otimes
\cV^{K3}_5.
\end{equation}
Factorisation homology over $E$ with the 24 punctures as marked
insertions produces the K3 chiral bialgebra
$\mathbf H_{\Delta_5}$:
\begin{equation}
\label{eq:ch-ce-k3-avatar-homology}
\int_E \cF^{\mathrm{ch}}_{\fg_{\Delta_5}, \cV^{K3}_5; E\setminus\{p_i\}}
\;\otimes\; \prod_{i=1}^{24}[\mathrm{insertion\,at\,}p_i]
\;=\;
\mathbf H_{\Delta_5}.
\end{equation}
\end{theorem}

\begin{proof}
Substitution of $C = E, \fg = \fg_{\Delta_5}, M = \cV^{K3}_5$ into
Definition~\ref{def:ch-ce-cochains}
gives~\eqref{eq:ch-ce-k3-avatar-factalg}. The 24 punctures parametrise
the 24 M5-brane insertions of
Definition~\ref{def:hcs-k3xE}, localised at the Kodaira $I_1$ fibres of
a nodal elliptic surface $E^{\mathrm{nod}}_{24}$ (equivalently, the 24
points of $E$ where the K3 elliptic fibration degenerates; Gaiotto
$\Sigma_{0,24}$ data). The factorisation-homology integral combines
the Dolbeault Fubini on $E\setminus\{p_i\}$ with the
Sym-part-CE closure, producing the K3 chiral bialgebra
$\mathbf H_{\Delta_5}$ by
Theorem~\ref{thm:g-delta5-chiral-ce-factalg-E} of
\texttt{celestial\_holography.tex}. The 24 insertions provide the
Jacobi-theta denominator $\eta^{24}$ on the partition function, which
by Borcherds' theta-correspondence gives the full $1/\Phi_{10}$
generating function.
\end{proof}

\begin{remark}[Concordance with Vol III Stage-2 factorisation]
\label{rem:ch-ce-vol3-concordance}
\index{Vol III!two-stage factorisation concordance}%
Theorem~\ref{thm:ch-ce-k3-avatar-24-punctures} is the Stage-2
specialisation of the Vol~III two-stage factorisation
$\Phi_d = \mathrm{Sp}^{\mathrm{ch}}_{\Sigma_{d-1},C}\circ\Phi^{\mathrm{FA}}_d$
at $d = 2$ (K3 fibre) with reference curve $C = E$ and
$\Sigma_{d-1} = \mathrm{K3}$. The Stage-1 $\Phi^{\mathrm{FA}}_2$
produces the $E_2$-holomorphic factorisation algebra on K3; Stage-2
$\mathrm{Sp}^{\mathrm{ch}}_{\mathrm{K3},E}$ specialises to the
$E_1$-chiral shadow on $E$ via $\int_{\mathrm{K3}}$. The chiral CE
factorisation algebra $\cF^{\mathrm{ch}}_{\fg_{\Delta_5}, \cV^{K3}_5;
E}$ is the explicit chain-level realisation of the Stage-2 output on
$E$.
\end{remark}

%%%%%%%%%%%%%%%%%%%%%%%%%%%%%%%%%%%%%%%%%%%%%%%%%%%%%%%%%%%%%%%%%%%%%%%%%%%%%
\section{Fifteen generator-relation identities}
\label{sec:ch-ce-fifteen-identities}
%%%%%%%%%%%%%%%%%%%%%%%%%%%%%%%%%%%%%%%%%%%%%%%%%%%%%%%%%%%%%%%%%%%%%%%%%%%%%

The following fifteen identities form the core presentation of
$\cF^{\mathrm{ch}}_{\fg,M;C}$ as a $\Disk_1$-algebra with explicit
generators and relations. Each identity is tagged by its claim-status
and cross-referenced to the corresponding proposition/theorem/proof.

\begin{enumerate}[label=\textup{(I\arabic*)},ref=I\arabic*]
\item\label{id:01}
$d_{\mathrm{ch}}\cA^a = \dbar_C\cA^a + f^a_{bc}\normord{\cB^b\cA^c}$
(gauge-field BRST; Proposition~\ref{prop:ch-ce-omega-linearity}\,(iii),
matching the BV-BRST field of
Definition~\ref{def:hcs-bv-brst}).
\ClaimStatusProvedHere
\item\label{id:02}
$d_{\mathrm{ch}}\cB^a = \dbar_C\cB^a + \tfrac12 f^a_{bc}\cB^b\cB^c$
(antighost BRST; Proposition~\ref{prop:ch-ce-omega-linearity}\,(iii)).
\ClaimStatusProvedHere
\item\label{id:03}
$\{\dbar_C, d_{\CE}\} = 0$ on
$\Omega^{0,\bullet}(C)\otimes\Sym^\bullet(\fg^\vee[-1])\otimes M$
(tensor-orthogonality; Proposition~\ref{prop:ch-ce-nilpotency}).
\ClaimStatusProvedHere
\item\label{id:04}
$\cA^a(z)\cA^b(w)\sim\delta^{ab}/(z-w)^2 + f^{ab}_c\cA^c(w)/(z-w)
+ \mathrm{regular}$
(tree-level OPE; Proposition~\ref{prop:classical-MC-relations}).
\ClaimStatusProvedHere
\item\label{id:05}
$\normord{\cA^a\cA^b}(z)\cA^c(w)\sim\mathsf C_{\mathrm{ad}}(\fg)
f^{ab}_d f^{dc}_e\cA^e(w)/(z-w)^2 + (\zeta(2)/(2\pi i)^2) f^{ab}_c
\partial\cA^c(w)/(z-w)$
(one-loop bubble; Theorem~\ref{thm:ch-ce-oneloop-motivic}).
\ClaimStatusProvedHere
\item\label{id:06}
$I^{\mathrm{sun}}_{\fsl_2}(z,w) = -\tfrac18\cdot f^{ab}_d f^{cd}_e
f^{ea}_f\cA^f(w)/(z-w)^3 + \mathrm{subleading}$
(two-loop sunrise coefficient; Theorem~\ref{thm:ch-ce-sunrise-sl2-minus-eighth}).
\ClaimStatusProvedHere
\item\label{id:07}
$\cF_{\fg,M}(D_1\sqcup D_2)\simeq\cF_{\fg,M}(D_1)\otimes\cF_{\fg,M}(D_2)$
(tensor-product axiom on disjoint discs;
Theorem~\ref{thm:ch-ce-tensor-axiom}).
\ClaimStatusProvedHere
\item\label{id:08}
$\cF_{\fg,M}(V) = \mathrm{hocolim}_i\cF_{\fg,M}(U_i)$ for Weiss cover
$\{U_i\}$ of $V$
(Weiss-cover descent; Theorem~\ref{thm:ch-ce-weiss-descent}).
\ClaimStatusProvedHere
\item\label{id:09}
$\Disk_3^Y = \Disk_1^{\R}\otimes_D\Disk_1^{\C}\otimes_D\Disk_1^{\C}$
(Dunn decomposition; Theorem~\ref{thm:ch-ce-disk3-algebra}).
\ClaimStatusProvedElsewhere
\cite[Thm.\,4.4.1]{CG17}
\item\label{id:10}
$\int_{\R^3\times\C^2}\Fobs^q_{\hCS}\simeq\cF^{\mathrm{ch}}_{\fg,M;C}$
(factorisation-homology Fubini; Theorem~\ref{thm:ch-ce-6d-hcs-reduction}).
\ClaimStatusProvedHere
\item\label{id:11}
$\cF^{\mathrm{ch}}_{\fg_{\Delta_5},\cV^{K3}_5; E\setminus\{p_i\}}$,
with $24$ punctures and factorisation-homology integral
$\int_E\cdot\prod[p_i]$, equals
$\mathbf H_{\Delta_5}$
(K3 avatar; Theorem~\ref{thm:ch-ce-k3-avatar-24-punctures}).
\ClaimStatusProvedHere
\item\label{id:12}
$\cA^a(z)\cdot W^V_\gamma = \int_\gamma\rho(T^a)W^V_{[\gamma_{\le z}]}
W^V_{[\gamma_{>z}]}\,dz$
(Wilson-line insertion; Theorem~\ref{thm:ch-ce-wilson-line-module}).
\ClaimStatusProvedHere
\item\label{id:13}
$T(z) = \tfrac{1}{2(k+h^\vee)}\sum_a\normord{J^aJ^a}(z)$
(Sugawara conformal vector at non-critical level;
Theorem~\ref{thm:ch-ce-topologization-E3-top}).
\ClaimStatusProvedElsewhere
\cite[Prop.\,5.1]{FrenkelBenZvi2004}
\item\label{id:14}
$\{\dbar_C, L_{-1}\}\omega = \omega$ on $\Omega^{0,\bullet}(D_z)$
(topologisation chain-homotopy witness;
Theorem~\ref{thm:ch-ce-topologization-E3-top}).
\ClaimStatusProvedHere
\item\label{id:15}
$\hbar^n\cdot K^\kappa_n = (-1)^n$ for the universal $\mathsf{B}$-family
at loop order $n$ (Bruinier--Heegner reciprocity;
Remark~\ref{rem:ch-ce-transcendence-bound}).
\ClaimStatusProvedElsewhere
\cite{Bruinier2002}
\end{enumerate}

\begin{remark}[Status distribution]
\label{rem:ch-ce-status-distribution}
\index{claim-status distribution}%
Of the fifteen core identities:
\begin{itemize}
\item \ClaimStatusProvedHere: identities~\eqref{id:01}--\eqref{id:08},
\eqref{id:10}, \eqref{id:11}, \eqref{id:12}, \eqref{id:14}
(eleven identities, with proofs in this chapter or referenced proofs in
Chapter~\ref{ch:six-d-hcs-e3-chiral-avatar-platonic}).
\item \ClaimStatusProvedElsewhere: identities~\eqref{id:09},
\eqref{id:13}, \eqref{id:15} (four identities, proved in
\cite{CG17,FrenkelBenZvi2004,Bruinier2002}).
\item \ClaimStatusConjectured: none.
\end{itemize}
\end{remark}

%%%%%%%%%%%%%%%%%%%%%%%%%%%%%%%%%%%%%%%%%%%%%%%%%%%%%%%%%%%%%%%%%%%%%%%%%%%%%
\section{Primary-literature concordance}
\label{sec:ch-ce-primary-lit-concordance}
%%%%%%%%%%%%%%%%%%%%%%%%%%%%%%%%%%%%%%%%%%%%%%%%%%%%%%%%%%%%%%%%%%%%%%%%%%%%%

\begin{remark}[Side-by-side correspondences]
\label{rem:ch-ce-side-by-side}
\index{primary literature!side-by-side}%
Each construction of this chapter has a direct primary-literature
counterpart:
\begin{center}
\begin{tabular}{ll}
\toprule
Construction here & Primary source \\
\midrule
$C^\bullet_{\mathrm{ch}}(\fg,M)$ tensor decomp.
& Francis--Gaitsgory 2012, Thm~6.4.2~\cite{FG12}\\
$\Disk_3$-action on $Y = \R^3\times C\times\C^2$
& Costello--Gwilliam Vol.~I, Thm~4.4.1~\cite{CG17}\\
Factorisation-homology Fubini
& Ayala--Francis--Tanaka 2017, Thm~4.1.3~\cite{AyalaFrancisTanaka2017}\\
One-loop Feynman Wick contraction
& Costello 2011, \S9--10~\cite{CostelloRenormalization}\\
Motivic-MZV basis at weight $n+1$
& Brown 2014, Thm~3.1~\cite{Brown2014}\\
6d hCS $E_3$ gen-rel
& Costello--Francis--Gwilliam 2026, \S3~\cite{CFG2026}\\
Bruinier--Heegner reciprocity
& Bruinier 2002, Prop~5.1~\cite{Bruinier2002}\\
Sugawara conformal vector
& Frenkel--Ben-Zvi 2004, Prop~5.1~\cite{FrenkelBenZvi2004}\\
Arnold relation on $\Conf_n(C)$
& Arnold 1969~\cite{Arnold1969}; Brieskorn 1973\\
\bottomrule
\end{tabular}
\end{center}
Each row is an equivalence at the level of chain complexes or of
$(\infty,1)$-operads. The cross-volume formula and object consistency
is verified against Vol.~I
\texttt{chapters/examples/landscape\_census.tex} and Vol.~III
\texttt{chapters/examples/cy\_d\_kappa\_stratification.tex}.
\end{remark}

\begin{remark}[Synthesis]
\label{rem:ch-ce-synthesis}
\index{chiral CE factorisation algebra!synthesis}%
The chiral CE factorisation algebra $\cF^{\mathrm{ch}}_{\fg,M;C}$ is
the chain-level realisation of the Stage-2 specialisation of the
Vol.~III CY-to-chiral functor $\Phi_d$ at $d = 2$, carrying fifteen
explicit generator-relation identities with motivic-MZV coefficients
and $\Disk_3$-structure via Dunn decomposition. Topologisation at
non-critical level produces $E_3^{\mathrm{top}}$ at each generic
bulk point, completing the $\mathsf{SC}^{\mathrm{ch,top}}$-to-$E_3$
promotion of Chapter~\ref{V1-chap:climax-platonic} of Vol.~I. The
K3 avatar on $C = E$ with 24 punctures produces the K3 chiral bialgebra
$\mathbf H_{\Delta_5}$, establishing a chain-level equivalence between
the 6d holomorphic Chern--Simons theory on
$\R^3\times\mathrm{K3}\times\C^2$ and the Borcherds--Kac--Moody
algebra $\fg_{\Delta_5}$ of the Gritsenko--Nikulin $\Phi_{10}$
denominator.
\end{remark}
